\documentclass[11pt]{diazessay} % Font size (can be 10pt, 11pt or 12pt)
\usepackage{ctex}
\usepackage{url}%引用网页
\title{\textbf{遗传经济学原理读后感} \\ {\Large\itshape Principles of Genetic Economics}} % Title and subtitle

\author{\textbf{穆思成} \\ \textit{XAUAT}} % Author and institution

\date{\today} % Date, use \date{} for no date

%中英文摘要
\newcommand{\enabstractname}{Abstract}
\newcommand{\cnabstractname}{摘要}
\newenvironment{enabstract}{%
\par\small
\noindent\mbox{}\hfill{\bfseries\enabstractname}\hfill\mbox{}\par
\vskip 2.5ex}{\par\vskip2.5ex}
\newenvironment{cnabstract}{
	\par\small
	\noindent\mbox{}\hfill{\bfseries\cnabstractname}\hfill\mbox{}\par
	\vskip 2.5ex}{\par\vskip2.5ex}

\renewcommand\cnabstractname{\textbf{摘要}}

\begin{document}

\maketitle % Print the title section

\begin{enabstract}
Genetic economics is an emerging interdisciplinary field that deeply integrates economic principles with genetic theory, particularly genomic data, 
to explore the biological underpinnings of human economic behavior at the molecular level. Founded by several scholars in 2007, 
it not only revolutionizes economics but also bridges multiple disciplines. 
Technological advancements in the 21st century, such as the substantial reduction in DNA sequencing costs and developments in computer technology,
have provided essential data and tools for genetic economics. The discipline employs \textbf{three core methodologies} – genome-wide association studies (GWAS),
Mendelian randomization, and direct application of genetic data in economic models – to uncover the genetic foundations of individual behavior, 
thereby furnishing a scientific basis for policy formulation.

The integration of Wen Tiejun's \textbf{"de-Western discourse construction"} emphasizes a non-Western perspective within genetic economics, 
highlighting localized research. This focus fosters a multiunderstanding of economic behavior and promotes global scientific dialogue and cooperation. 
With the rise of genetic economics, this text also combines insights from \textbf{"The Crowd: A Study of the Popular Mind"} 
illustrating the genetic determinants of individual behavior and offering a new materialistic perspective to economic research.

\end{enabstract}

\hspace*{3.6mm}\textit{Keywords:} three core methodologies，"de-Western discourse construction", "The Crowd: A Study of the Popular Mind"% Keywords

\vspace{80pt} 

\begin{cnabstract}

遗传经济学是将经济学原理与遗传学理论，特别是基因数据深度融合的新兴交叉学科，旨在从分子层面探索人类经济行为的生物学根源。
该领域由多位学者于2007年开创，它不仅革新了经济学，还在多学科间搭建桥梁。
21世纪的技术进步，如DNA测序成本的大幅降低和计算机技术的发展，为遗传经济学提供了必要的数据和工具。
该学科通过全基因组关联分析、孟德尔随机化法和基因数据在经济模型中的直接应用这\textbf{三大方法}，揭示了个体行为的遗传基础，为政策制定提供科学依据。
本文还结合温铁军的\textbf{“去西方话语主义建构”}，遗传经济学强调非西方视角，关注本土化研究，推动了对经济行为的多元化理解，并促进全球科学对话与合作。
随着遗传经济学的兴起，本文还结合了\textbf{《乌合之众》}的群体心理分析，展现了个体行为的遗传决定因素，为经济学研究提供了物质层面的新视角。

\end{cnabstract}
\hspace*{3.6mm}\textit{关键词:} 三大方法,“去西方话语主义建构”,《乌合之众》  % Keywords

\vspace{30pt} % Vertical whitespace between the abstract and first section

\begin{figure}[htbp]
    \centering
    \includegraphics[width=1\linewidth]{2}
    \caption{摘要图}
	\label{2}
\end{figure} 

\section*{引言}
读完《遗传经济学原理》这本书并查阅了相关资料后，首先惊奇于中国开拓遗传学理论与经典经济学结合工作的领导者朱晨老师，
竟然是如此年轻漂亮的一个学者，并在关注了朱老师的微信公众号和知乎账号后感受到了，老师对于扩大遗传经济学研究影响力的决心。
让我想起了2023年付汉良老师与其他老师一同开展的神经工程管理的第一次研讨会，各位老师都在共同努力去推进自己的研究，
通过各种方式（包括开展我们本次的通识课），努力向更多人传播生物神经遗传领域与各领域结合的研究成果。
这种科学探索的精神值得我们以及每一个青年学生去学习、借鉴、传承。

遗传经济学，由康奈尔大学经济学家Daniel Benjamin（现在南加大）、
哈佛大学经济学家Edward Glaeser，冰岛大学心血管遗传学家Vilmundur Gudnason、
联合学院心理学学家Christopher Chabris等学者于2007年创立\cite{1}。
作为一门新兴的交叉学科，不仅在经济学领域内掀起了波澜，更在生物学、心理学、遗传学乃至社会学诸多领域间架起桥梁，
它以独特的视角重新审视人类行为决策的根源——基因。

所有人都说21世纪是生物的世纪，我认为没错，在人工智能高速发展的今天我们可以利用计算机进行一系列我们之前没有过对于生物的解读学习。
结合我自己的体验，我在神经网络的学习中发现，现在的神经网络就是利用\textbf{分治法}，将一个大的问题拆分成海量小问题，
再从信息处理角度对\textbf{人脑神经元}进行抽象， 建立某种简单模型，按不同的连接方式组成不同的网络之后解决无数个小问题，最后得到解决大问题的方法，
这个过程也有黑盒效应，因为人类现在还不能理解神经元解决海量问题的原理和方法，抽象提取的特征值无法被解释，
只能按照仿生的逻辑输入，再得到结果，随着生物科技和计算机网络的发展，我相信将会有越来越多的成果呈现在世人面前。

遗传经济学，作为一门跨越传统经济学与生物学的桥梁，它在21世纪初露面世，由一群富有前瞻性的学者如Benjamin、Glaesler、Gudason、Chabris等联手开启。
这门学科的诞生，不仅是对经济学领域的一次革新，更是对人类自身认知的深潜心探索的里程碑。
它将经济学原理与现代遗传学的理论，特别是基因数据的深度整合，从分子层面挖掘个体行为的生物学根源，以全新的维度理解社会经济行为的决策机制。
本文将遗传经济学与《乌合之众》\cite{2}一书以及温铁军的“去西方话语主义建构”\cite{3}的看法相结合，尝试从不同的学科视角去理解遗传经济学，以期对遗传经济学的发展和应用提供新的视角和思路。

但是为什么是现在这个特殊的时间点将遗传与经济学相结合，这与计算机、生物科学、统计学的发展密不可分。

在20世纪90年代，随着DNA测序技术的发展，遗传学领域取得了巨大的进步，
人类基因组计划、人类遗传变异数据库、人类基因组变异图谱等数据集的建立，
为遗传学领域的研究提供了丰富的数据资源。而且随着新一代测序技术的出现，基因测序成本大幅降低，
个人全基因组测序成本迅速以超“摩尔定律（Moore’s Law）”的速度下降到了2017年的1120美元（见图\ref{1}NIH数据）\cite{1}，
且耗时大大缩短，大大推动了基因检测技术和遗传数据在各个领域中的应用。

\begin{figure}[htbp]
    \centering
    \includegraphics[width=1\linewidth]{1}
    \caption{NIH数据}
	\label{1}
\end{figure} 
同时，随着计算机技术的发展，遗传学领域开始探索利用计算机算法来处理和分析遗传数据，
如神经网络、遗传算法、基因表达分析、群体遗传学等。
这些技术的发展，为遗传学领域的研究提供了新的工具和手段。

\section{遗传经济学的三驾马车}

行为经济学作为经典经济学向非理性维度的探索，揭示了人们决策背后的心理与社会动因。而遗传经济学则更进一步，
它架起连接生物学基础与经济行为研究的桥梁，深入到遗传物质层次去剖析人的选择与行动模式。该领域的核心探索可归纳为三驾马车\cite{5}：

首先，遗传经济学采纳全基因组关联分析（GWAS）技术，这一前沿手段遍历人类基因组，旨在识别与经济行为特质紧密相关的遗传标记。
GWAS擅长追踪那些微妙的基因序列变化，它们虽小却能显著影响人的性格倾向，比如亲社会行为或冲动控制。
这些影响既包括直接由基因编码的“近基因性状”，也覆盖那些受基因间接作用、易受外界条件调制的“远基因性状”，如教育成就和风险承受能力。

其次，学科运用了孟德尔随机化法（MR），巧妙地将基因作为自然实验的工具变量，以绕过传统因果推断中的混杂因素。基因的随机分配特性，
如同自然界的随机对照试验，让经济学家得以从遗传的视角，无偏见地探究诸如健康状况或消费行为背后的真正驱动因素，这为揭示因果关系开辟了新的途径。

最后，遗传经济学直接整合基因数据进入经济学模型，作为先天条件的代表，这不仅清晰地区分了基因与环境对个体经济行为的不同贡献，
还为精确测量环境影响提供了可能。基因变量的这种直接应用，无论是作为研究焦点还是作为控制变量，
都极大增强了模型解释现实世界经济现象的能力，同时也促进了对后天因素影响的精细解析。

遗传经济学的发展，正如基因组学揭示的那样，仅是揭开人类行为奥秘的序幕。它不仅响应了现代社会对个体行为多维度理解的需求，
也与古老的东方智慧“天人合一”不谋而合，强调自然规律与人类活动的和谐统一。这一交叉学科的兴起，
预示着在构建具有中国特色的哲学社会科学体系中，遗传经济学将扮演一个不可或缺的角色，引领我们更深入地理解经济行为的生物学根源及其社会意义。

\section{结合温铁军的“去西方话语主义建构”分析遗传经济学应用基础}
整个研究是基于新古典经济学和行为遗传学，将研究扩展至基因层面。它认为，经济行为不仅受后天环境因素影响，还受遗传因素决定，强调基因在决策中的作用。
书中第七章给出了一系列在人类进化中基因随着环境的变化不断筛选的例子，讲述了人类发展过程中基因的几次大的筛选，以时间线，文明演化进程为纵向，以人类生活方式，
人类某些特定基因组为横向，交叉阐述了人类进化的客观历史进程与基因自然选择进程，最后形成了现代的人类，并且还说明了因为生产方式的快速巨大变革，人类的“自然选择”变成另一种方式，
但是仍然存在，在基因方面的选择依然存在，所以证明了研究的必要以及可实施性。

\textbf{为什么我要结合温铁军在“去西方话语主义建构”的看法？}
温老从唯物史观出发，立足客观经验，阐述了全球南方国家的客观历史进程。
他从新鲜经验和前沿学术成果发出，重看人类万年文明史，揭示人类文明发展的决定性因素，解释西方崛起与中国衰落的根本原因。
他的核心观点\textbf{人类多元文明的派生性差异化源于地理气候条件}与朱晨老师的自然选择基因筛选的观点并行，我认为如果可以放在一起解释将会有更立体的发现。

\begin{wrapfigure}{l}{6cm} % 靠文字内容的左侧
\includegraphics[width=6cm]{3}
\caption{去西方话语主义建构}
\label{3}
\end{wrapfigure}

遗传经济学的兴起，打破了西方经济学长期受新古典和新古典经济学对“理性人”假设的依赖，转而向生物遗传决定论。
温老的“去西方”思想促使我们审视遗传经济学，不局限于西方经济行为研究，而从非西方文化、社会经济背景，关注中国等非西方国家的基因-环境交互作用。
让我们的视角从局限的西方经济学，转向全人类的社会环境、文明演变的研究，
朱晨老师作为中国遗传经济学的第一人，也是帮助我们国家拥有大历史观，为构建人类命运共同体提供科学依据的学者，有助于我们重塑世界话语体系。

遗传经济学在“去西方”视角下，强调本土问题的探索。
温铁军鼓励本土化研究，中国作为5000年历史的东方文明，拥有着丰富的历史遗产，
其\textbf{现代经济行为的遗传经济学研究应考虑基因-环境因素}，如文化、教育、营养认知、食物多样性等，反映本土化研究，为政策提供科学依据。
在遗传经济学的研究中，我们应以本土数据为基础，以本土视角进行研究，制定更符合中国历史，中国国情的政策。

遗传经济学的跨学科性与温铁军的“去西方”相辅相得益。
遗传经济学要求经济学与其他学科如遗传学、心理学、流行病学、社会学、遗传学等深度整合，、
这与温铁军提倡的跨学科交流、知识融合相一致，推动对经济行为的多元理解，创新研究方法和理论。

“去西方”强调非西方数据的重要性。\textbf{遗传经济学应用中，数据收集本土基因数据库是关键}。
中国经济学界需要构建大规模的基因数据库，中国有着全世界最详细的历史资料，特别是清明两朝300多年间的县志，记录的最为详细，是中国研究经济学发展的独有优势。
温老有着30年下乡的乡土经验，看过无数的真实案例，这些案例有些已经被销毁，
如何尽可能得到更加全面真实的数据，将是提高研究的针对性的关键，
结合中国本土基因数据与历史事实数据，进行客观分析是未来遗传经济学的研究重点。
只有这样才能为政策制定提供科学依据，如农村发展、营养改善、学生健康政策等。

温铁军“去西方”视角下，遗传经济学能增强国际对话。
2023年以来，朱晨作为遗传经济学方向负责人，与中国医学科学院、哈佛大学等全球顶尖科研机构的学者们组建了跨学科团队，
重点关注社会隔离、生活环境等社会经济因素对中国和欧美居民健康和疾病发生的因果性影响，
前期研究成果已发表在《柳叶刀》子刊《eClinicalMedicine》（共同第一作者，影响因子15.1），引起国内外公众和媒体的广泛关注，并有一系列后续研究工作推进中\cite{4}。
通过遗传经济学研究，中国学者在全球遗传学、经济学领域内建立话语体系，如SSGA、国际合作，反哺”西方中心主义，提升中国学术地位，促进全球科学对话与合作。
以及近10几年来提出的丝绸之路经济带和人类命运共同体，给世界提供了之前从没有过的基于亚洲大陆一体化的全球建设经验。

\section{结合《乌合之众》一书阐述基因数据的经济学价值}

《乌合之众》一书，作为一部关于群体心理的经典著作，描绘了群体行为与思维的非理性特质，揭示了个体在集体中的心理动态。
古斯塔夫·勒庞的观点用一句话概括就是：他认为\textbf{群体是愚蠢的}。
遗传经济学，正是将群体行为与思维的非理性特质用基因的方法解释，揭示个体在集体中的心理动态，
将其视角转向遗传经济学领域，特别是基因数据的价值，我发现遗传经济学利用个体的遗传信息为经济学提供了前所未有的洞见地揭示行为决策机制。
\begin{wrapfigure}{l}{5cm} % 靠文字内容的左侧
	\includegraphics[width=5cm]{4}
	\caption{《乌合之众》}
	\label{4}
	\end{wrapfigure}
如同《乌合之众》揭示群体行为的非理性，遗传经济学通过基因数据深入理解个体行为的差异。
遗传信息帮助经济学家区分个体的异质性，即个体特质差异不仅源于后天（遗传），更深入揭示了个体行为的物质基础，与乌合之众的内在一致性与差异性。

遗传经济学强调基因与环境交互，如同乌合之众群体与环境的关系，基因数据能揭示环境如何影响个体。
它提供了一种理解基因与社会经济行为的物质层面如何在个体差异性状与远基因性状（如教育、收入）与环境交互，如同乌合之众中的个体在社会环境中的位置，影响。
\textbf{这里的环境是指的社会环境，就是个体在群里里的环境，与上一章节讨论的群体在客观自然环境中的环境不同，本文正是从宏观到微观的对遗传经济学的阐释。}

遗传数据的随机性，孟德尔化方法如同乌合之众的随机性，为因果推断提供了自然实验的替代。
基因数据的随机性，使经济学家能够利用孟德尔随机化作为工具变量，如同乌合之众的群体随机性，控制环境变量，为因果推断提供了科学的工具。

乌合之众的群体决策影响，基因数据提供政策影响。遗传经济学通过群体特征的经济策略制定，如教育、健康、福利、收入分配等。
遗传数据指导政策，如乌合之众的群体决策影响，为政策制定提供了科学依据，如同乌合之众的群体影响，使政策制定更科学、合理。

遗传经济学的基因数据如同乌合之众的反哺育科学。
遗传学，经济学的研究为遗传学提供了新视角，基因-表型机制，通过行为研究，如同乌合之众的反哺育了遗传学的内在机制研究，为遗传学提供了经济学的视角，如同乌合之众的反馈。

因此，基因数据在遗传经济学的运用如同《乌合之众》的群体心理，揭示了个体行为的内在机制，为经济学提供了物质层面的视角，理解行为与决策的科学工具，
揭示了个体的物质基础，为经济学研究提供了科学依据，利用遗传学的视角我们可以分析股票市场的变化，很多时候股票市场会随着市场情绪的变化而变化，
\textbf{市场情绪就如同乌合之众般的人们内心深处的共同情绪变化}，如能深度挖掘人们内心的遗传性行为，
有利于投资金融市场的宏观调控，为政策提供了科学的决策，促进了经济学与遗传学的互哺，如同乌合之众的社会科学的桥梁。

\section{问题与回答}
\textbf{我的问题：}
在读到第10章水稻种植与现代人基因这部分时，让我产生了一个疑问就是中国地域辽阔，既有稻作农业也有麦作农业，
也有这两种种植方式交融的地区，而且因为一些历史原因人类大规模迁移的情况不在少数，单纯的讨论遗传基因与两种耕作方式的关系是否能够成立？
我认为气候变化是对于耕作方式以及人类迁徙影响根本原因，是不是应该从气候的角度讨论，而不是从耕作方式讨论。

所以我就此问题在知乎上询问了朱老师，并且很荣幸的获得了答复。
\begin{figure}[htbp]
    \centering
    \includegraphics[width=1\linewidth]{5}
    \caption{问题}
	\label{5}
\end{figure} 

\textbf{老师的回答：}
你好，确实有这种可能。我们在这篇论文中控制了经纬度等和温度相关的变量，
同时对现代化假设、传染病假设等可能性进行了检验，但受限于数据没有直接检验气候变化的影响。但你说的可能性是存在的。
如果未来有更好更合适的数据可以检验。我们现在的结果只能说是在当前可用数据、方法论和证据的条件下，所能得到的最可能的解释。
科学本身是一门螺旋发展的学科，未来新数据新方法表明目前的理解可能是不确定的、不完整的，甚至是错误的，都是在意料之中的。

\begin{figure}[htbp]
    \centering
    \includegraphics[width=1\linewidth]{6}
    \caption{回答}
	\label{6}
\end{figure} 

\vspace{30pt}

最后，真的感谢朱老师能够在网上认真回答我的问题，以及付老师以及各位老师8节课的指导，让我能够了解到神经工程管理以及遗传经济学的相关知识，
并且向我展示了科学工作者研究科学的过程，以及科研工作者包容且积极的精神。

\bibliographystyle{unsrt}

\bibliography{sample.bib}

\end{document}
